
This work builds on three research threads: dataset documentation completeness analysis, metadata standardization and conflicts, and FAIR data principles operationalization for ML.

\subsubsection{2.1 Dataset Documentation Completeness}

Yang et al. (2024) conducted large-scale analysis of dataset card completion on HuggingFace, examining 7,433 datasets and revealing marked heterogeneity in completion rates. Practitioners prioritize Dataset Description and Structure sections over Considerations sections, with completion correlating with dataset popularity. This work established that automated subsection-level completion analysis is feasible at scale, but focused on a single platform without examining cross-platform variation or explicitly measuring platform UX features.

This work extends Yang's analysis to cross-platform comparison across 3 ML repositories (OpenML, HuggingFace, UCI) with friction score as explanatory variable. Where Yang observed heterogeneity patterns, we measure UX-driven effects with quantitative comparisons.

Reid and Williams (2023) characterized voice dataset documentation practices through 13 machine learning practitioner interviews and rubric analysis of 9 datasets, finding that fragmented codification hinders comparison and combination across platforms. Their qualitative analysis highlighted standardization needs but did not quantify platform-level effects.

\subsubsection{2.2 Metadata Conflicts and Standardization}

Strecker (2026) investigated metadata conflicts across 8 geoscience and social science repositories, identifying both implementation conflicts (workflows, decisions) and inter-standard conflicts (inherent schema differences) contributing to incomplete DataCite metadata. This work established that metadata completeness is multifaceted, driven by both technical and conceptual factors.

This work builds on Strecker's taxonomy by measuring outcomes (presence rates) rather than conflict sources, linking platform UX features to completeness patterns at 10,000+ dataset scale in the ML domain. Where Strecker provided qualitative classification, we provide quantitative effect sizes (45-61pp for optional fields).

Batzner et al. (2026) developed a shared schema for AI evaluation results, successfully ingesting 22,235+ model evaluations from heterogeneous sources via automated converters. This work demonstrated that friction reduction (automated converters from popular formats) enables large-scale voluntary contribution.

Huang et al. (2025) systematically assessed metadata completeness in the Gene Expression Omnibus (GEO) repository, analyzing 253 studies with 164,000 samples and finding 25\% critical metadata omitted with only 11.5\% complete phenotype sharing. Public repositories contained 62\% phenotypes (3.$5\times$ more than publications alone), with non-human samples showing better metadata completeness than human studies.

Kim et al. (2025) proposed tier-based standards for FAIR sequence data sharing in microbiome research, analyzing 2,929 publications and finding nearly half do not meet minimum data availability standards. Poor metadata standardization creates barriers to harmonization.

\subsubsection{2.3 FAIR Data Principles and ML Workflows}

Samuel et al. (2020) applied FAIR data practices to ML pipelines, proposing provenance capture using ProvBook tool with Jupyter Notebooks for end-to-end reproducibility. They demonstrated that factors beyond source code and datasets influence reproducibility.

Logan et al. (2023) reviewed machine learning datasets in mammography for FAIR principles adherence, finding variability in interoperability and dataset skew toward clinical use-cases. They recommended improving interoperability through BIRADS criteria adherence and consistent file formats.

Giner-Miguelez et al. (2025) analyzed 4,041 scientific data papers for ML-requested dimensions, comparing coverage and trends with NeurIPS Datasets and Benchmarks venue. They provided recommendation guidelines for data creators and publishers.

\subsubsection{2.4 Positioning}

Prior work established that (1) dataset documentation is heterogeneous \citep{Yang2024Navigating}, (2) metadata conflicts stem from implementation and schema differences \citep{Strecker2026Metadata}, (3) automated infrastructure enables contribution \citep{Batzner2026Every}, and (4) FAIR practices matter for reproducibility (Samuel 2020, Logan 2023). However, no study quantitatively linked repository UX design to metadata completeness outcomes across multiple platforms.

This work introduces friction score as a quantifiable UX metric, measures cross-platform presence patterns at 10,000+ dataset scale, and validates enforcement versus friction-reduction mechanisms.
